Matching generation time improves the predefined fresh 48-step chunk relative to applying its old delay offset: the paired mean gain is 0.579 (Table~\ref{tab:decoder}). This establishes an execution-index effect for that fixed option in the new cohort. Its time-matched fresh versus stale difference remains unresolved, as does its descriptive advantage over continued local control ($+0.120$, 95\% interval $[-0.187,0.433]$). Improvement in this decoder comparison does not by itself establish a useful refresh channel or an overall advice advantage.
\begin{table}[tb]\centering\small
\caption{Independent generation-time decoding control: four primary reward contrasts on 60 new environments. Time uses generation age; legacy uses request delay. Stale decoding is identical for both. Approximate 98.75\% stratified paired-bootstrap intervals retain the four-contrast family. Fixed48 applies the predefined source suffix on every environment; the frozen gate uses unchanged choices across fresh decoders. This study does not alter the original channel estimates.}\label{tab:decoder}
\setlength{\tabcolsep}{3pt}\begin{tabular}{llrr}\toprule Receiver & Contrast & Mean & Interval\\\midrule
Fixed48 & Fresh: time $-$ legacy & +0.579 & [+0.173, +0.997]\\
Fixed48 & Time-fresh $-$ stale & +0.149 & [-0.204, +0.516]\\
Frozen $P,F$ & Fresh: time $-$ legacy & -0.030 & [-0.298, +0.192]\\
Frozen $P,F$ & Time-fresh $-$ stale & -0.188 & [-0.541, +0.116]\\
\bottomrule\end{tabular}\end{table}


The frozen $P,F$ gate responds differently. It chooses a suffix in 24/60 environments (21 for 48 steps and three for 16), so time matching changes 24 action sequences while changing no fresh-channel choice. Its mean decoder effect is $-0.030$, with a primary interval crossing zero. Local-option selections are unaffected. The fixed option's overall gain therefore does not transfer automatically to the gate's selected states and intervention lengths. The gate still predicts legacy-execution gains; this study does not test learning or recalibrating it for the new rule. Complete channel/decoder outcomes, success endpoints and index checks are in Online Resource 1.

These new results qualify a practical reading of the earlier channel test: literal replacement in the legacy receiver and generation-matched execution answer different questions. We retain both, with separate environments, receivers and inference families. The positive fixed-chunk decoder result is not used to relabel the original refresh effect or the packet-independent $+0.735$ counterexample.
